\begin{ouverture}[college2]{10}{Exploiter la gravitation universelle}{Partie 2 · Mouvements et interactions}{Comment comparer l’attraction entre deux objets ?}
\objectif{Utiliser la loi de gravitation lorsqu’elle est fournie.}
\objectif{Convertir distances et masses en unités SI.}
\objectif{Distinguer gravitation, poids et impesanteur.}
\prerequis{Puissances de dix et calculatrice.}
\prerequis{Forces à distance ; P = mg.}
\end{ouverture}
\pcPage{Le cours}{Exploiter la gravitation universelle}
\phantomsection\label{ess-10}
\essSection{1}{Une interaction attractive}
Deux objets possédant une masse s’attirent par gravitation. Les forces sont dirigées selon la droite joignant leurs centres dans le modèle de corps sphériques. Chaque objet exerce une force sur l’autre. Les valeurs sont égales et les sens opposés ; les forces s’appliquent à des objets différents. Ce modèle permet notamment d’étudier l’interaction Terre-Lune.\par
\pcFigure{figures/b10-forces.pdf}{Deux forces sur deux systèmes ; tailles et distance ne sont pas à l’échelle.}
\essSection{2}{Utiliser la relation fournie}
Quand l’énoncé la donne, on exploite la relation \[F=G\dfrac{m_A m_B}{d^2}.\] F est en N, les masses en kg et d en m, distance entre les centres. On emploie par exemple $G=6{,}67\times10^{-11}\unite{N\,m^2/kg^2}$. Écrire les conversions avant le calcul, regrouper le dénominateur et vérifier la touche « puissance » de la calculatrice.\par
\clearpage
\essSection{3}{Comparer sans tout recalculer}
À distance inchangée, doubler une masse double F ; doubler les deux masses multiplie F par quatre. À masses inchangées, doubler d divise F par quatre ; tripler d divise F par neuf. Le carré porte sur toute la distance. Ces comparaisons concernent les conditions annoncées et le même modèle.\par
\pcFigure{figures/b10-distance.pdf}{La force diminue comme l’inverse du carré de la distance, pour des masses fixées.}
\essSection{4}{Du modèle au poids}
Près de la surface d’un astre, son attraction domine et correspond au poids dans le modèle simplifié. Avec un rayon R, on peut comparer $P=mg$ et $F=GMm/R^2$, sans confondre rayon et altitude. Un satellite en orbite reste attiré : l’impesanteur ressentie à bord correspond à une chute commune, et non à une disparition de la gravitation.\par
\begin{pcEncadre}{method}{Un exemple expliqué}Deux masses de 10 kg et 20 kg, distantes de 2 m, sont modélisées par $F=Gm_Am_B/d^2$. $F=6{,}67\times10^{-11}\times10\times20/2^2=3{,}335\times10^{-9}\unite N$. Cette force est très faible à l’échelle d’un dynamomètre scolaire ; elle n’est pas nulle.\end{pcEncadre}
\pcPage{Expérimenter et retenir}{Exploiter la gravitation universelle}
\phantomsection\label{activite-10}
\begin{pcEncadre}{method}{Investigation en classe}Utiliser une simulation ou un tableur, avec la loi fournie. Fixer les masses et calculer F pour d, 2d et 3d ; puis fixer d et changer une seule masse. Tracer des points et expliquer pourquoi F en fonction de d ne forme pas une droite. Aucun montage de gravitation sensible n’est exigé au collège.\end{pcEncadre}
\pcSousTitre{Ce que je dois savoir expliquer}\begin{questions}
\item Identifier d comme distance entre les centres.
\item Calculer avec des kg, m et une loi fournie.
\item Justifier les facteurs 2, 4 ou 9 par la relation.
\end{questions}
\begin{pcEncadre}{attention}{Une erreur à éviter}Une altitude de 400 km n’est pas la distance au centre de la Terre. Dans ce modèle, il faut ajouter le rayon terrestre avant de convertir et de mettre au carré.\end{pcEncadre}
\pcRepere{cor-activite-10}{Pistes pour l’investigation}\pcGrille
\clearpage\markright{Je démarre}
\rubrique{Je démarre}
\par\noindent\begin{minipage}{\linewidth}
\begin{exercice}[1]{Nature de l’action}
\begin{questions}
\item La gravitation est-elle une action de contact ?
\item Est-elle attractive dans le modèle étudié ?
\end{questions}
\end{exercice}
\end{minipage}
\par\noindent\begin{minipage}{\linewidth}
\begin{exercice}[1]{Deux objets}
\begin{questions}
\item La Terre attire la Lune. La Lune attire-t-elle la Terre ?
\item Sur quel objet agit la force exercée par la Lune ?
\end{questions}
\end{exercice}
\end{minipage}
\par\noindent\begin{minipage}{\linewidth}
\begin{exercice}[1]{Unités}
\begin{questions}
\item Dans $F=Gm_A m_B/d^2$, donner les unités de m et d pour obtenir F en N.
\item Convertir 6 400 km en m.
\end{questions}
\end{exercice}
\end{minipage}
\par\noindent\begin{minipage}{\linewidth}
\begin{exercice}[1]{Quelle distance ?}
\begin{questions}
\item Pour deux sphères, quelle distance figure dans la loi donnée ?
\item Pourquoi la distance entre leurs surfaces ne suffit-elle pas ?
\end{questions}
\end{exercice}
\end{minipage}
\par\noindent\begin{minipage}{\linewidth}
\begin{exercice}[1]{Une masse doublée}
\begin{questions}
\item La masse A double, le reste ne change pas. Que devient F ?
\item Si les deux masses doublent ?
\end{questions}
\end{exercice}
\end{minipage}
\clearpage\markright{Je m’entraîne}
\rubrique{Je m’entraîne}
\par\noindent\begin{minipage}{\linewidth}
\begin{exercice}[2]{Une distance doublée}
\begin{questions}
\item Les masses sont fixées et d double. Que devient F ?
\item Et si d est divisée par deux ?
\end{questions}
\end{exercice}
\end{minipage}
\par\noindent\begin{minipage}{\linewidth}
\begin{exercice}[2]{Une force modèle}
\begin{questions}
\item Avec $G=6{,}67\times10^{-11}$ en unités SI, calculer F pour 1 kg, 1 kg et 1 m.
\item Est-ce une force nulle ?
\end{questions}
\end{exercice}
\end{minipage}
\par\noindent\begin{minipage}{\linewidth}
\begin{exercice}[2]{Une comparaison rapide}
\begin{questions}
\item À d, F vaut 36 N. Quelle valeur à 3d avec les mêmes masses ?
\item Quelle valeur à 2d ?
\end{questions}
\end{exercice}
\end{minipage}
\par\noindent\begin{minipage}{\linewidth}
\begin{exercice}[2]{Depuis l’orbite}
\begin{questions}
\item Rayon terrestre donné : 6 400 km ; altitude : 400 km. Déterminer d en km puis m.
\item Quel terme doit-on mettre au carré ?
\end{questions}
\end{exercice}
\end{minipage}
\par\noindent\begin{minipage}{\linewidth}
\begin{exercice}[2]{Un petit calcul}
\begin{questions}
\item Calculer F pour $m_A=100\unite{kg}$, $m_B=50\unite{kg}$, $d=10\unite m$, $G=6{,}67\times10^{-11}$ en SI.
\item Prévoir la valeur si seule $m_B$ devient 100 kg.
\end{questions}
\end{exercice}
\end{minipage}
\clearpage\markright{Je raisonne}
\rubrique{Je raisonne}
\par\noindent\begin{minipage}{\linewidth}
\begin{exercice}[3]{Trois facteurs}
\begin{questions}
\item La masse A double et d double, B reste fixe. Par quel facteur F est-elle multipliée ?
\item Justifier avec la relation.
\end{questions}
\end{exercice}
\end{minipage}
\par\noindent\begin{minipage}{\linewidth}
\begin{exercice}[3]{Un faux équilibre}
\begin{questions}
\item « Les deux forces Terre-Lune s’annulent, donc la Lune n’est pas attirée. » Quelle erreur de système apparaît ?
\item Quelle force faut-il retenir pour étudier seulement la Lune ?
\end{questions}
\end{exercice}
\end{minipage}
\par\noindent\begin{minipage}{\linewidth}
\begin{exercice}[3]{Impesanteur}
\begin{questions}
\item Pourquoi l’impesanteur dans une station orbitale ne signifie-t-elle pas absence d’attraction ?
\item Quelle hypothèse donne l’explication qualitative ?
\end{questions}
\end{exercice}
\end{minipage}
\par\noindent\begin{minipage}{\linewidth}
\begin{exercice}[3]{Retrouver g}
\begin{questions}
\item La loi fournie donne une force de 588 N sur une masse de 60 kg près du sol. Trouver g à partir de P = mg.
\item Quelle masse aurait un poids de 98 N au même lieu ?
\end{questions}
\end{exercice}
\end{minipage}
\par\noindent\begin{minipage}{\linewidth}
\begin{exercice}[3]{La précision du modèle}
\begin{questions}
\item Une calculatrice affiche dix chiffres pour F, mais les distances ont deux chiffres significatifs. Pourquoi éviter une conclusion à dix chiffres ?
\item Citer une hypothèse du modèle gravitationnel utilisé.
\end{questions}
\end{exercice}
\end{minipage}
\pcPage{Défi du chapitre}{Exploiter la gravitation universelle}
\phantomsection\label{defi-ex-10}
\pcDefiTitre{10}{Autour d’une planète}
\begin{questions}
\item Une planète a $M=6{,}0\times10^{24}\unite{kg}$. Un satellite de 500 kg est à $d=7{,}0\times10^6\unite m$. Calculer F avec $G=6{,}67\times10^{-11}$ en SI.
\item Si d double, que devient F ?
\item Comparer les forces planète-satellite et expliquer pourquoi leurs effets ne sont pas identiques.
\end{questions}
\pcGrille
